% ============================================================================
% PART II: PHYSICS ORGANIZED BY OBSERVABLE FAMILY
% Section 6: Topological Sectors in 3+1 Dimensions
% ============================================================================
%
% Purpose
%   The four-dimensional gauge-theory consequences of topology: theta vacua,
%   topological susceptibility chi_t, U(1)_A problem, eta' mass, strong CP,
%   neutron EDM, axion physics. The narrative chain is standard in the axion
%   / strong-CP literature (Hook, Dine, Marsh):
%
%     topology -> vacuum structure -> chi_t -> E(theta) -> eta' mass;
%     neutron EDM constrains bar-theta; axion dynamically cancels it.
%
%   This is the "spectroscopic / CP-sensitive" observable family.
%
% Status: NEW WRITING. Build from the outline in
%   tqft_observables_literature_review.md Sec. 5, supplemented with any class
%   notes available in PROJECTS/QFT/literature/ for the instanton / theta
%   material.
%
% Length target: ~15--20 pages
%
% Subsection plan (from Final QFT Project Overview.md)
%   6.1  Theta vacua, topological susceptibility, and the U(1)_A problem
%   6.2  The eta' mass and the Witten-Veneziano formula
%   6.3  Strong CP, the neutron EDM, and the axion
%   6.4  When the observables are clean, and when they are only effective
%
% References
%   't Hooft 1976 BellJackiw (\cite{tHooft1976BellJackiw})
%   't Hooft 1976 Pseudoparticle (\cite{tHooft1976Pseudoparticle})
%   Witten 1979 (\cite{Witten1979U1GoldstoneBoson})   -- U(1)_A half of WV
%   Veneziano 1979 (\cite{Veneziano1979U1WithoutInstantons}) -- WV other half
%   Di Vecchia-Veneziano 1980 (\cite{DiVecchiaVeneziano1980ChiralDynamics})
%     -- large-N chiral lagrangian with theta-dependence
%   Dine 2000 (\cite{Dine2000TASIStrongCP})            -- strong CP TASI
%   Hook 2018 (\cite{Hook2018TASIStrongCPAxions})      -- modern TASI review
%   Marsh 2016 (\cite{Marsh2016AxionCosmology})        -- axion phenomenology
%   Teper 2000 (\cite{Teper2000TopologyInQCD})         -- lattice topology review
%   Cichy et al. 2015 (\cite{CichyLatticeWittenVeneziano})
%     -- lattice test of Witten-Veneziano
%   Arean et al. 2021 (\cite{AreanU1AHolographicQCD})  -- holographic chi_t
%   Abel et al. 2020 (\cite{Abel2020NeutronEDM})       -- nEDM world limit
%   Experiment placeholders: ADMX, CAST/IAXO, CASPEr, MADMAX
%     (\cite{ADMXOverviewPlaceholder}, \cite{CASTIAXOOverviewPlaceholder},
%      \cite{CASPErOverviewPlaceholder}, \cite{MADMAXOverviewPlaceholder})
%
% Writing notes
%   - Open by reminding the reader that Sec. 3 already introduced the theta
%     term and instanton number. Here we follow the chain from topology to
%     measurable observables.
%   - 6.1 derives chi_t as the curvature of E(theta).
%   - 6.2 states Witten-Veneziano and motivates large-N; cite lattice support.
%   - 6.3 connects bar-theta -> nEDM and describes axion solution (PQ
%     mechanism, axion mass from chi_t, coupling hierarchy by UV model).
%   - 6.4 is a local "iffiness" subsection paralleling Sec. 5.6.

\section{Topological Sectors in $3+1$ Dimensions}
\label{sec:sectors-3plus1d}

% TODO: Open with the reminder that theta in F wedge F-tilde is not just an
% observable --- it is a coupling whose physical consequences run through the
% entire hadronic and CP-violation story.

\subsection{$\theta$ Vacua, Topological Susceptibility, and $U(1)_A$}
\label{subsec:theta-chi-u1a}

% TODO: Define chi_t = int d^4x <q(x) q(0)> with q = (g^2/32 pi^2) tr(F F-tilde).
% Relate to E(theta) = E(0) + (1/2) chi_t theta^2 + O(theta^4) at small theta.
% Explain U(1)_A: anomaly kills naive Goldstone; large-N saves the story.
% Cite: \cite{tHooft1976BellJackiw}, \cite{tHooft1976Pseudoparticle},
%       \cite{Witten1979U1GoldstoneBoson}, \cite{Veneziano1979U1WithoutInstantons},
%       \cite{DiVecchiaVeneziano1980ChiralDynamics},
%       \cite{Teper2000TopologyInQCD}

\subsection{The $\eta'$ Mass and the Witten--Veneziano Formula}
\label{subsec:eta-prime-wv}

% TODO: Derive (sketch) the Witten-Veneziano relation
% m_{eta'}^2 + m_{eta}^2 - 2 m_K^2 = (2 N_f / F_pi^2) chi_t^{YM}
% and interpret chi_t^{YM} as a pure-gauge-theory topological observable
% inferred from hadron spectroscopy.
% Cite: \cite{Witten1979U1GoldstoneBoson}, \cite{Veneziano1979U1WithoutInstantons},
%       \cite{DiVecchiaVeneziano1980ChiralDynamics},
%       \cite{CichyLatticeWittenVeneziano}, \cite{AreanU1AHolographicQCD}

\subsection{Strong CP, the Neutron EDM, and the Axion}
\label{subsec:strong-cp-axion}

% TODO: bar-theta = theta + arg det M_q; nEDM bound d_n <= 1.8e-26 e cm ->
% |bar-theta| <~ 10^{-10}. PQ mechanism: promote bar-theta to dynamical axion
% field a/f_a; axion potential V(a) from chi_t gives m_a^2 f_a^2 = chi_t;
% small-angle relaxation dynamically enforces CP.
% Phenomenology: axion mass-coupling plane, haloscope / helioscope /
% CASPEr / MADMAX experimental classes.
% Cite: \cite{Dine2000TASIStrongCP}, \cite{Hook2018TASIStrongCPAxions},
%       \cite{Marsh2016AxionCosmology}, \cite{Abel2020NeutronEDM},
%       \cite{ADMXOverviewPlaceholder}, \cite{CASTIAXOOverviewPlaceholder},
%       \cite{CASPErOverviewPlaceholder}, \cite{MADMAXOverviewPlaceholder}

\subsection{When the Observables Are Clean, and When They Are Only Effective}
\label{subsec:qcd-caveats}

% Dedicated "iffiness" subsection for the 3+1d story
% (tqft_observables_literature_review.md Sec. 5.4):
%   - F wedge F-tilde is locally a total derivative, but not globally trivial.
%   - Instanton-gas reasoning is semiclassical, not the strong-coupling story.
%   - chi_t is Euclidean and theory-side; lattice / holographic QCD are needed
%     to connect to hadronic observables.
%   - Neutron EDM measurement is clean; extraction of bar-theta from the
%     measurement involves hadronic matrix elements and model dependence.
%   - Axion phenomenology: the mass comes from chi_t and is robust; photon /
%     nucleon / spin couplings depend on the UV completion (KSVZ vs DFSZ etc).
% Cite: \cite{Hook2018TASIStrongCPAxions}, \cite{Marsh2016AxionCosmology},
%       \cite{Teper2000TopologyInQCD}
